\documentclass[11pt, a4paper]{article}
\usepackage[margin=1in]{geometry}
\usepackage{amsmath}
\usepackage{amsfonts}
\usepackage{amsthm}
\usepackage{tikz}
\usepackage{xcolor}
\usepackage{framed}
%\theoremstyle{definition}
\newtheorem{definition}{Definition}[section]
\newtheorem{theorem}{Theorem}[section]
\newtheorem{corollary}{Corollary}[section]

\title{Everything Summer `26}
\author{Absur Khan Siam}
\date{\today}

\begin{document}
\maketitle
\section{Kitaev Quantum Double Model}
The flow: Fix a lattice $\rightarrow$ edges with quantum degrees of freedom $\rightarrow$ local operators and Hamiltonian $\rightarrow$ ground state and excitations.

\textit{State space:} Lattice with V, E, and F: the set of vertices, edges, and faces, respectively. Assign an orientation to each $e \in \mathrm{E}$ (arbitrary).
\begin{equation*}
    \text{Sites } \mathcal{S} = \{(p,~v)~|~p \in F,~v \in \partial p\}.
\end{equation*}
At every edge $e \in \mathrm{E}$, we put a $|G|$-dimensional qudit $\mathcal{H}_e = \mathbb{C}[G]$ with computational basis $\{|g\rangle\}$. Total Hilbert space $\mathcal{H}_\text{tot} = \otimes_{e \in E} \mathcal{H}_e$ with $\text{dim}(\mathcal{H}_\text{tot}) = |G|^{|E|}$. Also,
\begin{equation*}
    |g_{e_1}~g_{e_2} \dots g_{e_{|E|}}\rangle := |g_{e_1}\rangle \otimes |g_{e_1}\rangle \otimes \dots \otimes |g_{e_{|E|}}\rangle.
\end{equation*}
Note: The ring multiplication structure of $\mathbb{C}[G]$ comes into play in the local operators.

\textit{Hamiltonian:} $A_g(v)$ is defined as follows: \textcolor{red}{(Figure to be included)}
\begin{equation*}
    |\dots g_1~g_2~g_3~g_4\dots\rangle \mapsto |\dots gg_1~gg_2~g_3g^{-1}~g_4g^{-1}\dots\rangle.
\end{equation*}
(Note that $A_{g_1}A_{g_2} = A_{g_1g_2}$ since $g_2^{-1}g_1^{-1} = (g_1g_2)^{-1}$.)
We define the vertex operator $A(s)$ where $s \in \mathcal{S}$:
\begin{equation*}
    A(s) := A(v) = \frac{1}{|G|}\sum_g A_g(v)
\end{equation*}
This is a projector. \textit{Proof}: $A^2(s) = \frac{1}{|G|^2}\sum_{g_1,~g_2}A_{g_1}A_{g_2}(v) = \frac{1}{|G|^2}\sum_{g_1,~g_2}A_{g_1g_2}(v)$.\\
Now, due to the rearrangement of finite groups, $\frac{1}{|G|^2}\sum_{g_1,~g_2}A_{g_1g_2}(v) = \frac{1}{|G|^2}\sum_{g_1}\left(\sum_gA_{g}(v)\right) = \frac{1}{|G|^2}\cdot|G|\cdot \sum_g A_g(v) = A(v)$. \qed

$B_h(v)$ is defined as follows: \textcolor{red}{(Figure to be included)}

If $h = e$, the relation $g^{o(e_1)}g^{o(e_2)}g^{o(e_3)}g^{o(e_4)} = e$ can be cyclically permuted ($o(e)$ is the orientation of edge $e$). Hence, the starting vertex for traversing the plaquette $p$ \textbf{does not} matter.

$A_g(s)$ and $B_h(s)$ do not commute but satisfy \begin{equation}
    A_g(s)B_h(s) = B_{ghg^{-1}}(s)A_g(s).
\end{equation}

In a more convenient notation, $A_g(v)B_h(p) = B_{ghg^{-1}}(p)A_g(v)$. This means that we start traversing the plaquette $p$ counterclockwise from vertex $v$. \textcolor{red}{(Figure to be included)}

Firstly,
\begin{align*}
    A_g(v)B_h(p)|\dots g_1~g_2~g_3~g_4\dots\rangle
    &= \delta_{h,~g_2^{-1}g_3g_4g_1^{-1}}A_g(v)|\dots g_1~g_2~g_3~g_4\dots\rangle\\
    &= \delta_{h,~g_2^{-1}g_3g_4g_1^{-1}}A_g(v)|\dots gg_1~gg_2~g_3g^{-1}~g_4g^{-1}\dots\rangle.
\end{align*}
Here, we ignore the effects of $A_g(v)$ on the rest of the vertices since they are irrelevant in this context. Secondly,
\begin{align*}
    & B_{ghg^{-1}}(p)A_g(v)|\dots g_1~g_2~g_3~g_4\dots\rangle\\
    &= B_{ghg^{-1}}(p)|\dots gg_1~gg_2~g_3g^{-1}~g_4g^{-1}\dots\rangle\\
    &= \delta_{ghg^{-1},~\left(g_2g^{-1}\right)^{-1}g_3g_4\left(gg_1\right)^{-1}}A_g(v)|\dots gg_1~gg_2~g_3g^{-1}~g_4g^{-1}\dots\rangle.
\end{align*}
Finally, we note that the conditions
\begin{align*}
    ghg^{-1} &= \left(g_2g^{-1}\right)^{-1}g_3g_4\left(gg_1\right)^{-1}\\
    \Leftrightarrow ghg^{-1} &= g\left(g_2^{-1}g_3g_4g_1^{-1}\right)g^{-1}\\
    \Leftrightarrow h &= g_2^{-1}g_3g_4g_1^{-1}
\end{align*}
are all equivalent, meaning that
\begin{equation*}
    A_g(v)B_h(p) = B_{ghg^{-1}}(p)A_g(v) \qed
\end{equation*}

{\footnotesize If we flip the direction of $e_1$, we get $h = g_2^{-1}g_3g_4g_1$ and $ghg^{-1} = \left(g_2g^{-1}\right)^{-1}g_3g_4\left(g_1g^{-1}\right)$, which are also equivalent. The argument is similar if we flip $e_2$.}

They satisfy the same relations as the generators of $\mathrm{DG}$, which implies that the algebra of local operators at a (?) site $s = (p, v)$ generated by $A_g$ and $B_h$ form an $N = |G|^{|E|}$ (the exponent $|E|$, for `a' site, how?) dimensional matrix representation of the quantum double algebra $DG$ given by the homomorphism
\begin{align*}
    \rho_s: \mathrm{DG} &\to GL(N, \mathbb{C})\\
    D_{e,~g} &\mapsto A_g(v)\\
    D_{h,~e} &\mapsto B_h(p).\\
\end{align*}
Altogether, the quantum double \textbf{Hamiltonian} is given by
\begin{equation}
    H = -\sum_{v \in V} A(v) - \sum_{p \in P}B(p),
\end{equation}
where $B(p) := B_e(p)$. We can check from the non-commuting relation of $A_g$ and $B_h$ that all the projectors in the Hamiltonian commute with each other.

\end{document}